\documentclass[aspectratio=169]{beamer}
\usetheme{metropolis}
\metroset{sectionpage=none, numbering=counter, progressbar=frametitle}
\usepackage{amsmath, amssymb}
\usepackage[T1]{fontenc}

\title{Domain-Partition Algorithm for Turbomachinery Blades}
\subtitle{Slide 9: Postprocessing (Periodic Boundaries)}
\date{\today}
\author{}
\institute{}

\begin{document}

\begin{frame}{Postprocessing --- Periodic Boundary Closure}

\begin{itemize}
    \item \textbf{Collapse 3-sided seam wedges} (Xiao-style simplification)
    \item \textbf{Master-slave seam conformity} (DLR Sauer/Morsbach 2023 sec 2.7)
    \item \textbf{Result:} periodic block structure with conforming seam
\end{itemize}

\vspace{0.5em}
\small
\textit{Code:} \texttt{collapse\_seam\_wedges}, \texttt{symmetrize\_seam\_junctions} in \texttt{tmesh\_partition.py:332--488} \quad
\textit{Lit:} Sauer/Morsbach 2023, ``Automated blocking for turbomachinery'', sec 2.7

\note{
Speaker notes:
Two postprocessing steps enforce periodic boundary closure on the unwrapped hub cylinder.
First, collapse_seam_wedges (lines 332--402) removes 3-sided seam wedges: when a singularity near the theta seam sends two adjacent arms into the seam wall, the shorter arm is dropped and the two wall segments are re-joined at the freed junction. This is a Xiao-style simplification that eliminates degenerate triangles at the periodic boundary.
Second, symmetrize_seam_junctions (lines 405--488) implements master-slave conformity following DLR Sauer and Morsbach 2023 section 2.7. The left wall is the master; every junction is projected onto it, the right wall is rebuilt as the exact plus-pitch translate (copy plus shift, never re-projected), and the endpoints of interior curves are moved onto the exact junction coordinates. Mirrored junctions without an interior curve remain as hanging T-nodes. The result is a perfectly periodic block structure with a conforming seam.
}

\end{frame}

\end{document}
